\subsection{ENSO groups}
The ONI record (Figure~\ref{fig:oni}) yields \nSuper{} super El Niño winters (\superWinters), \nElnino{} other El Niño winters, \nNeutral{} neutral winters and \nLanina{} La Niña winters between 1950/51 and \lastWinter. Peak ONI values were 2.14 (1983), 2.37 (1998) and 2.59\,\degC{} (2016); the next strongest events were 2023/24 (1.99), 1972/73 (1.84) and 1965/66 (1.82).

\subsection{December to February}
\textbf{Temperature.} Detrended DJF mean temperature in super El Niño winters exceeded neutral winters in every mainland Nordic country (Figure~\ref{fig:tmean}, Table~\ref{tab:contr}): Finland \fiTSvNdiff\,\degC{} (\fiTSvNlo{} to \fiTSvNhi, permutation $p = \fiTSvNp$), Sweden \seTSvNdiff\,\degC{} (\seTSvNlo{} to \seTSvNhi, $p = \seTSvNp$), Norway \noTSvNdiff\,\degC{} (\noTSvNlo{} to \noTSvNhi, $p = \noTSvNp$), Denmark \dkTSvNdiff\,\degC{} (\dkTSvNlo{} to \dkTSvNhi, $p = \dkTSvNp$) and the Nordic mean \nordicTSvNdiff\,\degC{} (\nordicTSvNlo{} to \nordicTSvNhi, $p = \nordicTSvNp$). Iceland went the other way (\isTSvNdiff\,\degC, \isTSvNlo{} to \isTSvNhi). Ordinary El Niño winters were only marginally milder than neutral ones (Nordic \nordicTEvNdiff\,\degC, \nordicTEvNlo{} to \nordicTEvNhi, $p = \nordicTEvNp$), so most of the super-versus-neutral difference is also present as a super-versus-El Niño difference (Nordic \nordicTSvEdiff\,\degC, \nordicTSvElo{} to \nordicTSvEhi, $p = \nordicTSvEp$; Finland \fiTSvEdiff\,\degC, \fiTSvElo{} to \fiTSvEhi, $p = \fiTSvEp$). La Niña winters did not differ from neutral winters in any country.

The permutation $p$ values are the ones to read. Welch's $t$ gives $p$ between 0.04 and 0.07 for Sweden, Norway and the Nordic mean, but it estimates the spread of the super group from three winters that happen to resemble each other; the permutation test asks how often three winters drawn at random from the record produce a mean this far from the rest, and the answer is one time in three to one time in seven. The intervals quoted above for the super group are Welch $t$-intervals and share that fragility.

\textbf{Precipitation.} Finland was the only country with a coherent precipitation surplus in super El Niño winters: \fiPSvNdiff\,mm over DJF relative to neutral winters (\fiPSvNlo{} to \fiPSvNhi, $p = \fiPSvNp$) and \fiPSvEdiff\,mm relative to ordinary El Niño winters (\fiPSvElo{} to \fiPSvEhi, $p = \fiPSvEp$), about a quarter of the seasonal total, with the surplus present in all three winters (detrended anomalies of $+16$, $+23$ and $+34$\,mm in 1983, 1998 and 2016). Norway showed a surplus of similar size with a far wider interval (\noPSvNdiff\,mm, \noPSvNlo{} to \noPSvNhi), Sweden and Denmark showed nothing (\sePSvNdiff{} and \dkPSvNdiff\,mm), and Iceland was drier on two winters of data. The Nordic mean was \nordicPSvNdiff\,mm (\nordicPSvNlo{} to \nordicPSvNhi). Ordinary El Niño winters did not differ from neutral winters in precipitation anywhere (Figure~\ref{fig:precip}).

\textbf{Snow, frost and thaw.} Mean snow depth was lower in super winters (Nordic \nordicSSvNdiff\,cm, \nordicSSvNlo{} to \nordicSSvNhi, $p = \nordicSSvNp$; Finland \fiSSvNdiff\,cm, \fiSSvNlo{} to \fiSSvNhi; Figure~\ref{fig:snow}). Thaw days were more frequent (Nordic \nordicThawSvNdiff{} days, \nordicThawSvNlo{} to \nordicThawSvNhi, $p = \nordicThawSvNp$; Finland \fiThawSvNdiff{} days, \fiThawSvNlo{} to \fiThawSvNhi, $p = \fiThawSvNp$; Figure~\ref{fig:thaw}) and frost days slightly fewer (Finland \fiFrostSvNdiff{} days, \fiFrostSvNlo{} to \fiFrostSvNhi). These are the same signal as the temperature difference seen through different thresholds, not independent evidence.

\textbf{Spatial pattern.} At station level (Figures~\ref{fig:map} to \ref{fig:heat2}) the warm anomaly is largest in southern Sweden, Denmark and southern Norway and fades northward; Icelandic stations are cool. The Finnish precipitation surplus is present at every FMI station with a complete record. Only a handful of the several hundred station-by-metric cells reach $p < 0.05$, consistent with chance in an uncorrected table.

\textbf{Absolute values.} The three super winters were not mild in absolute terms. National DJF means in Finland were $-6.5$, $-7.6$ and $-6.2$\,\degC{} in 1982/83, 1997/98 and 2015/16, that is anomalies of $+0.3$, $-0.6$ and $+0.6$\,\degC{} against 1991 to 2020 (Figure~\ref{fig:tsfi}, Supplementary Figure~S1 in \texttt{graphs\_export/}). They rank as mild only relative to the trend line of their own decade, and the detrended framing is what makes them look warm.

\subsection{January to March}
In JFM the temperature signal disappears (Table~\ref{tab:contrjfm}): Finland $-0.3$\,\degC, Sweden $+0.6$, Norway $+0.4$, Denmark $+1.0$, Iceland $-0.8$, Nordic mean $+0.2$\,\degC, all with intervals spanning zero and permutation $p \ge 0.3$ (Figure~\ref{fig:tmeanjfm}). The Finnish precipitation surplus shrinks to $+18$\,mm ($p = 0.26$). Ordinary El Niño winters are, if anything, slightly cooler than neutral in JFM in Finland ($-0.35$\,\degC), in the direction the late-winter stratospheric pathway predicts, but far from distinguishable from noise.

\subsection{Continuous ONI}
Across all labelled winters the DJF ONI had no monotonic relation with any national metric: Spearman's $\rho$ for DJF temperature was between $-0.02$ and $0.09$ and for precipitation between $-0.15$ and $0.07$ (Figure~\ref{fig:scatter}). Whatever the three super winters share, it is not the tail of a linear dose-response.
